\section{Testing Java}
\label{sec:testing}

\subsection{Struttura dei test}

Ho diviso i test in più classi, una per ogni aspetto che volevo coprire, così è più
facile capire cosa controlla cosa. Nella tabella qui sotto ho messo l'elenco con,
accanto, quello che ciascuna classe va a verificare.

\begin{longtable}{@{}p{4.3cm}p{8.7cm}@{}}
\toprule
\textbf{Classe} & \textbf{Copre} \\
\midrule
\endhead
\texttt{ParcheggioTest} & Costruttori, precondizioni violate, tutti i rami di \texttt{assegnaPosto}/\texttt{liberaPosto}, test parametrico CSV. \\
\texttt{SmartParkingFSMTest} & Tutti gli 8 stati e le transizioni, fallback disabile, parcheggio pieno, self-loop NEG/INGR/TARIF. \\
\texttt{McdcVerificaTest} & Copertura MC/DC delle due decisioni composte di \texttt{assegnaPosto} (Tabella~\ref{tab:mcdc1} e~\ref{tab:mcdc2}). \\
\texttt{DisplayMockTest} & Notifiche a \texttt{Display} verificate con Mockito. \\
\texttt{ScenarioAvallaTest} & Traduzione fedele di \texttt{test\_pieno.avalla} in JUnit (Model Based Testing). \\
\texttt{CombinatorialTest} & Tabella pairwise a 6 righe (Tabella~\ref{tab:pairwise}). \\
\texttt{UISeleniumTest} & Smoke test della UI (tag \texttt{ui}, escluso dalla build standard). \\
\bottomrule
\end{longtable}

\subsection{MC/DC}

Dentro \texttt{assegnaPosto} ci sono due decisioni composte, cioè con più condizioni
messe in and/or, e per queste ho voluto una copertura MC/DC vera e propria: per ogni
condizione serve una coppia di casi di test in cui, cambiando solo quella condizione,
cambia anche il risultato della decisione. Vediamole una alla volta.

La prima decisione è quella del ramo STD/ABBONATO: $D_1 = (\text{utente}{=}\text{STD}
\lor \text{utente}{=}\text{ABBONATO}) \land \text{postiStd}{>}0$.

\begin{longtable}{@{}p{1cm}p{0.6cm}p{0.6cm}p{0.6cm}p{0.6cm}p{2.6cm}p{5.5cm}@{}}
\toprule
Caso & A & B & C & $D_1$ & Esito & Coppia indipendente \\
\midrule
\endhead
TC1 & T & F & T & T & POSTO\_STD & base \\
TC2 & T & F & F & F & NESSUNO & isola C (vs TC1) \\
TC3 & F & F & T & F & POSTO\_DIS & isola A (vs TC1) \\
TC4 & F & T & T & T & POSTO\_STD & isola B (vs TC3) \\
\bottomrule
\caption{MC/DC -- Decisione 1.}
\label{tab:mcdc1}
\end{longtable}

La seconda decisione è quella del ramo DISABILE, cioè il fallback, con
$C_1{=}\text{postiDis}{>}0$ e $C_2{=}\text{postiStd}{>}0$.

\begin{longtable}{@{}p{1cm}p{1cm}p{1cm}p{2.6cm}p{6.5cm}@{}}
\toprule
Caso & $C_1$ & $C_2$ & Esito & Coppia indipendente \\
\midrule
\endhead
TC5 & T & --- & POSTO\_DIS & base \\
TC6 & F & T & POSTO\_STD & isola $C_1$ (vs TC5) \\
TC7 & F & F & NESSUNO & isola $C_2$ (vs TC6) \\
\bottomrule
\caption{MC/DC -- Decisione 2.}
\label{tab:mcdc2}
\end{longtable}

\subsection{Combinatorial testing (pairwise)}

L'assegnazione del posto dipende da tre parametri: il tipo di utente
(\texttt{tipoUtente} $\in$ \{STD, DISABILE, ABBONATO\}), i posti standard liberi
(\texttt{postiStd} $\in \{0, 1\}$) e quelli disabili liberi (\texttt{postiDis}
$\in \{0, 1\}$). Provare tutte le combinazioni sarebbe possibile ma inutile: mi basta
coprire tutte le \emph{coppie} di valori, ed è quello che fa la tabella pairwise qui
sotto, che con sole 6 righe le copre tutte.

\begin{longtable}{@{}p{0.6cm}p{2cm}p{1.6cm}p{1.6cm}p{3.5cm}@{}}
\toprule
\# & tipoUtente & postiStd & postiDis & Esito atteso \\
\midrule
\endhead
1 & STD      & 0 & 0 & NESSUNO \\
2 & STD      & 1 & 1 & POSTO\_STD \\
3 & DISABILE & 0 & 1 & POSTO\_DIS \\
4 & DISABILE & 1 & 0 & POSTO\_STD (fallback) \\
5 & ABBONATO & 0 & 1 & NESSUNO \\
6 & ABBONATO & 1 & 0 & POSTO\_STD \\
\bottomrule
\caption{Tabella pairwise (copre tutte le coppie dei 3 parametri).}
\label{tab:pairwise}
\end{longtable}

\subsection{Mockito}

La FSM, a ogni cambio di stato, manda un messaggio a un oggetto \texttt{Display}. Per
controllare che questi messaggi partano davvero, senza dover tirare in ballo la UI vera,
ho usato Mockito: creo un finto \texttt{Display} e poi con \texttt{verify} controllo che
il metodo \texttt{mostra} venga chiamato con il testo giusto. Nell'esempio qui sotto
riempio il parcheggio e verifico che, al terzo ingresso, il display riceva davvero
``Accesso negato''.

\begin{lstlisting}[language=Java, caption={Verifica delle notifiche con Mockito}]
@Test
void notificaAccessoNegatoQuandoParcheggioPieno() {
    SmartParkingFSM fsm = new SmartParkingFSM(new Parcheggio(0, 0), display);
    fsm.step(SensorInput.sensIn());
    fsm.step(SensorInput.utente(TipoUtente.STD));
    fsm.step(SensorInput.utente(TipoUtente.STD));
    verify(display).mostra("Accesso negato");
}
\end{lstlisting}

\subsection{Model Based Testing}

L'idea del model based testing è prendere uno scenario già scritto sul modello astratto
e tradurlo in un test sul codice vero. Ho preso lo scenario \texttt{test\_pieno.avalla} e
l'ho riportato passo passo in \texttt{ScenarioAvallaTest}: ogni blocco
\texttt{set}/\texttt{step}/\texttt{check} dell'avalla diventa la corrispondente riga
JUnit, e ho lasciato un commento accanto a ciascuna per richiamare la riga originale
dello scenario.

\subsection{Esecuzione e copertura}

\screenshot{junit_run_all_green}{Esecuzione di tutta la suite JUnit}{%
In Eclipse (o IntelliJ), tasto destro sulla cartella \texttt{src/test/java} $\to$
\textbf{Run As $\to$ JUnit Test}, oppure da terminale \texttt{mvn test}. Screenshot
della vista JUnit con tutti i test verdi e il conteggio totale (es. ``Runs: N/N,
Errors: 0, Failures: 0'').}

La copertura l'ho misurata con JaCoCo, lanciando \texttt{mvn verify} e guardando il
report in \texttt{target/site/jacoco/index.html}. Dal conteggio ho tolto la UI Vaadin
(che testo end-to-end con Selenium e che JaCoCo non riesce a misurare) e la classe di
avvio di Spring Boot, che sono solo 3 righe di boilerplate; l'esclusione è scritta nel
\texttt{pom.xml}. Alla prima misurazione il package \texttt{core} stava al 98\% di
instruction e al 95\% di branch. La classe \texttt{Parcheggio} era già al 100\%/100\%
(coerente con la verifica ESC dei contratti, \S\ref{sec:jml-ref}), mentre in
\texttt{SmartParkingFSM} restavano scoperti 3 branch: il self-loop di
\texttt{gestisciUsc} (\texttt{transito\_ok} falso in USC), il ramo
\texttt{display == null} di \texttt{cambiaStato} e il \texttt{default} dello switch di
\texttt{step}. I primi due erano veri buchi di test, che ho chiuso aggiungendo due
test mirati (\texttt{uscRestaInUscFinoAlTransito}, \texttt{fsmFunzionaSenzaDisplay}).
Il terzo invece non è raggiungibile: lo switch gestisce tutti e 8 i valori dell'enum
\texttt{StatoSistema}, per cui il \texttt{default} (difensivo, lancia
\texttt{IllegalStateException}) non può mai essere eseguito. È l'analogo Java dei rami
infattibili già osservati con AsmetaV sul modello (\S\ref{sec:asmetav-vc}). Dopo aver
aggiunto i due test, quel \texttt{default} resta l'unico branch non coperto del
package.

\screenshot{jacoco_coverage_report}{Report JaCoCo -- totale (solo package core)}{%
Pagina principale del report (stato finale, dopo l'aggiunta dei due test mirati):
98\% instruction / 98\% branch coverage; l'unico elemento non coperto e' 1 branch su
61 (il \texttt{default} difensivo dello switch, strutturalmente irraggiungibile).}

\screenshot{jacoco_detail_parcheggio}{Report JaCoCo -- dettaglio per classe del package core}{%
Vista del package \texttt{it.tvsw.smartparking.core}: \texttt{Parcheggio} 100\%/100\%,
enum e \texttt{SensorInput} 100\%, \texttt{SmartParkingFSM} 96\%/96\%.}

\screenshot{jacoco_detail_fsm}{Report JaCoCo -- dettaglio per metodo di SmartParkingFSM}{%
Vista per metodo di \texttt{SmartParkingFSM}: tutti i metodi al 100\% tranne
\texttt{step()}, il cui unico branch non coperto è il \texttt{default} difensivo
dello switch, strutturalmente irraggiungibile (discusso nel testo).}
